\chapter{宠物养成对战斗能力的映射}

前文已经分别讨论了宠物成长、转生、融合、蛋喂药、乱敏、命中、会心和骑宠伤害。真正的研究价值并不在于这些部分各自成立，而在于它们能否被接成一条从“养成输入”到“战斗输出”的完整链条。本章的任务正是把这条链条显式写出来：宠物的成长率如何转化为攻击、防御、敏捷；转生、融合和喂药如何改变这种转化；而这些变化又如何进一步进入先手、命中、会心和最终伤害。

\section{从 \vstd{} 到战斗属性的映射}

就宠物本体而言，成长率并不直接等于战斗属性，而是先经过出生随机、升级随机与显示公式，才形成玩家可见的攻击、防御、敏捷。若把高等级时的原始累计点向量记为
\begin{equation}
    \bm{r}=(V_r,S_r,T_r,D_r),
\end{equation}
则第 3 章已经给出显示映射：生命主要由 $V_r$ 主导，攻击主要由 $S_r$ 主导，防御主要由 $T_r$ 主导，敏捷则近乎由 $D_r$ 单独决定。因此，从边际影响上看，四项成长对战斗能力的主要路径可以概括为
\begin{equation}
    V \rightarrow HP,\qquad S \rightarrow ATK,\qquad T \rightarrow DEF,\qquad D \rightarrow SPD.
    \label{eq:growth-main-path}
\end{equation}

但式~\eqref{eq:growth-main-path} 只是主路径，而不是完整路径。因为显示公式中仍存在交叉项：$V$ 和 $D$ 会少量进入攻击、防御，$S$ 和 $T$ 也会通过总和与随机路径间接影响成长评分。因此，所谓“主 V 宠”“主 D 宠”在源码层面更准确的说法应是“某一维在主路径上占优势”，而不是“其它三维可以忽略”。

进入战斗后，这些显示属性又会继续被转换成固定属性和战斗属性。对宠物而言，最关键的桥梁变量是敏捷。宠物显示敏捷越高，通常越容易抬高 \code{WORKFIXDEX} 和后续的 \code{WORKQUICK}，于是更容易在乱敏、回避、会心和反击中占优。也正因为此，在很多战斗目标中，成长里的 $D$ 维不仅仅对应“面板上多几点敏”，而是沿着一条更长的链条影响先手权与概率判定。

\section{从转生、融合到战斗能力的长期收益}

转生、融合和蛋喂药的共同点，在于它们都不是直接修改某个战斗公式中的输入，而是先修改“进入成长流程之前的成长档”。这意味着它们的作用具有长期累积性。设原始成长档为 $\bm{g}$，转生后成长档为 $\bm{g}^{(new)}$，融合蛋基础成长为 $\bm{b}$，喂药并孵化后的成长档为 $\bm{h}$，则战斗能力的比较实质上是在比较
\begin{equation}
    \Psi\bigl(\bm{g}\bigr),\qquad \Psi\bigl(\bm{g}^{(new)}\bigr),\qquad \Psi\bigl(\bm{h}\bigr),
\end{equation}
其中 $\Psi(\cdot)$ 表示“成长随机 $\rightarrow$ 高等级显示属性 $\rightarrow$ 战斗固定属性 $\rightarrow$ 战斗结果”的复合映射。

转生的长期收益主要体现在两个方面。第一，它通过转生能力值和按权重重分配，整体改变了四维成长的起点；第二，转生后的新成长率进入的是“转生宠 Rank 表”，从而改变了后续升级随机的区间结构。因此，转生不是简单的“总和变大”，而是同时改变总和、方向和 Rank 解释规则。

融合和喂药的长期收益则更像一个约束优化问题。融合先通过主宠 $0.6$ 与副宠平均后 $0.4$ 生成蛋基础成长，喂药再经过单项折算、50 压缩、按比例回分配和 150 压缩，最后得到孵化后成长 $\bm{h}$。因此，融合链条的长期收益不只取决于“投入了多少好素材”，还取决于这些素材和药效在多次取整与上限之后还剩多少。这种离散损耗正是许多玩家直觉与实际结果不一致的原因。

\section{从蛋喂养到战斗目标的定向优化}

由于喂药和孵化会在进入普通成长流程之前确定一组新的成长档，因此对战斗目标的优化应当首先落在 $\bm{h}$ 上，而不是落在蛋基础成长 $\bm{b}$ 上。若把一个战斗目标写成效用函数 $U$，则蛋喂养问题可抽象为
\begin{equation}
    \max_{\text{喂药方案}} U\bigl(\Psi(\bm{h})\bigr).
\end{equation}
在实际工具中，为了降低求解复杂度，常先用一个线性代理目标
\begin{equation}
    U_0(\bm{h}) = w_V h_V + w_S h_S + w_T h_T + w_D h_D
    \label{eq:linear-hatch-target}
\end{equation}
来近似真正的战斗目标。

式~\eqref{eq:linear-hatch-target} 的价值在于，它把复杂的战斗诉求先投影回孵化后成长空间。例如，若目标偏向先手与回避，则 $w_D$ 应更高；若目标偏向骑宠近战爆发，则 $w_S$ 与 $w_D$ 通常应同时抬高；若目标偏向承伤与稳定性，则 $w_V$ 与 $w_T$ 更重要。等得到候选解之后，再把它们带回更完整的战斗模型中做二次比较。

这一思路之所以合理，是因为喂药阶段离战斗结果中间还隔着出生后随机和升级随机，直接对最终战斗结果做全局枚举往往代价过高；而对 $\bm{h}$ 做代理优化，则能把高维随机问题拆成“先确定方向，再在随机层验证”的两步问题。当前项目里的融合蛋页面，正是沿着这一思路把“底层整数规则”和“目标导向搜索”接起来的。


\section{宠物忠诚失控（nono）与动作改写}

对宠物而言，成长并不是只通过面板三围进入战斗。当前源码里，宠物忠诚还会通过 \code{CHAR_WORKFIXAI} 这一工作值，直接改变“玩家选了什么动作，宠物最终执行什么动作”。这一机制的核心函数是 \code{BATTLE_PetLoyalCheck()}，它在每只单位正式执行动作之前调用；因此，它不是一个结算后附带触发的小概率事件，而是位于“动作写入之后、动作分发之前”的主流程分支。

这一点对理解宠物培养很重要。若把“宠物成长 $\rightarrow$ 战斗输出”的链条只写成属性映射，那么会忽略一个更上游的问题：宠物是否真的会按玩家指定的技能出手。对低忠诚宠而言，战斗输出不只是某个技能公式的结果，而是“原指令 $\rightarrow$ 忠诚检查 $\rightarrow$ 实际动作”的复合结果。

\subsection{忠诚阈值与改写结果}

设当前回合读取到的忠诚值为
\begin{equation}
    a = \code{CHAR\_WORKFIXAI},
\end{equation}
并记 \code{RAND(1,100)} 的结果为 $R$。源码不是只设一个“是否 nono”的单阈值，而是按 $a$ 的区间，把宠物分到不同改写模式，如表~\ref{tab:nono-modes} 所示。

\begin{table}[H]
    \centering
    \caption{宠物忠诚失控（nono）的动作改写区间}
    \label{tab:nono-modes}
    \begin{tabularx}{\linewidth}{@{}>{\raggedright\arraybackslash}p{2.5cm} >{\raggedright\arraybackslash}p{2.7cm} X@{}}
        \toprule
        忠诚区间 $a$ & 触发条件 & 本回合结果 \\
        \midrule
        $a \ge 80$ & 恒不触发 & 完全按玩家原指令行动 \\
        $70 \le a < 80$ & $R < 10$ & 9\% 概率进入“乱目标” \\
        $60 \le a < 70$ & $R < 20$ & 19\% 概率进入“乱目标” \\
        $50 \le a < 60$ & $R < 35$ & 34\% 概率进入“乱目标” \\
        $40 \le a < 50$ & $R < 50$ & 49\% 概率进入“乱目标” \\
        $30 \le a < 40$ & $R < 70$ & 69\% 概率改放另一项宠技 \\
        $20 \le a < 30$ & $R < 70$ & 69\% 概率改放另一项宠技 \\
        $10 \le a < 20$ & $R < 80$ & 79\% 概率攻击主人，否则攻击敌人 \\
        $0 \le a < 10$ & $R < 60$ & 59\% 概率攻击主人，否则直接逃跑 \\
        \bottomrule
    \end{tabularx}
\end{table}

表~\ref{tab:nono-modes} 中的概率之所以是 $9\%$、$19\%$、$34\%$ 而不是看上去更直觉的 $10\%$、$20\%$、$35\%$，原因在于源码使用的是严格不等式 \code{<}，并且随机值来自闭区间整数 \code{RAND(1,100)}。例如“\code{R < 10}”只会在 $R=1,\dots,9$ 时成立。

这里还要强调，只有 $20 \le a < 40$ 这一段的“随机动作”模式，才会把玩家原本指定的宠技 A 改写成另一项宠技 O；$40 \le a < 80$ 只会乱目标，$0 \le a < 20$ 则会进一步退化成攻击主人、攻击敌人或逃跑。因此，玩家口中那种“明明点了 A，结果放成另一项技能 O”的 nono 现象，并不是所有低忠诚区间都会发生，而主要集中在 $20 \le a < 40$。

\subsection{为何会出现“技能 A 的面板修正保留，但实际放出技能 O”}

这一现象并不是玩家错觉，而是源码执行顺序的直接结果。宠物技能在玩家输入阶段就会先写入一组本回合工作值，例如
\begin{equation}
    \begin{aligned}
        (&\code{WORKBATTLECOM1},\code{WORKBATTLECOM2},\code{WORKBATTLECOM3},\\
         &\code{WORKATTACKPOWER},\code{WORKDEFENCEPOWER},\code{WORKQUICK},\dots).
    \end{aligned}
\end{equation}
随后在战斗主循环中，系统先做状态序列与基础修正，再调用 \code{BATTLE_PetLoyalCheck()} 检查忠诚；只有在这一步之后，才按 \code{WORKBATTLECOM1} 进入具体动作分支。于是，若宠物先按玩家原指令 A 写入了攻、防、敏等工作值，而忠诚检查又把动作编号改成另一项技能 O，那么是否出现“先吃到 A 的面板效果，再放出 O”，就取决于 O 会不会把这些工作值重新覆盖。

源码里这类情况确实存在。例如 \code{PETSKILL_AttackCrazed()} 会先把
\begin{equation}
    \code{WORKATTACKPOWER} \leftarrow 0.8 \, \code{WORKFIXSTR},
    \qquad
    \code{WORKDEFENCEPOWER} \leftarrow 0.7 \, \code{WORKFIXTOUGH},
\end{equation}
再写入 \code{BATTLE_COM_S_ATTCRAZED}；而 \code{PETSKILL_AttackShoot()} 在当前源码中只改写动作编号、目标和 \code{WORKBATTLECOM3}，其原本的攻防重写代码被整体注释掉了。于是，若宠物先选中了前者，但在忠诚失控的“随机动作”分支里又被改写成后者，就会出现“前一个技能留下的攻防修正仍在，最终放出的却是后一个技能”的连锁效果。

从研究方法上说，这说明宠物战斗输出不能只写成“技能 O 的伤害公式”。对低忠诚宠，更准确的表达应是
\begin{equation}
    \text{最终输出} = \Omega\bigl(\text{原指令 }A,\; \text{忠诚区间},\; \text{改写结果 }O,\; \text{O 是否覆盖 A 已写入的工作值}\bigr).
\end{equation}
也就是说，低忠诚宠的实战收益不仅由成长与技能池决定，还由“本回合动作写入顺序”和“各技能对工作槽的覆盖范围”共同决定。这也是为什么某些表面上只看技能文本无法解释的爆发现象，在源码里却能由忠诚改写链得到直接说明。

\section{不同培养目标下的宠物类型}

若把战斗目标按主导维度划分，可以把宠物大致分成四类。第一类是\emph{先手/控制型}，核心在于尽可能抬高 $D$，使宠物本体或骑宠混合后的速度敏更容易跨过乱敏阈值。第二类是\emph{近战爆发型}，更重视 $S$ 与 $D$ 的组合：$S$ 负责提高攻击路径，$D$ 负责抢先手和提高会心、命中相关概率。第三类是\emph{承伤稳定型}，其重点通常落在 $V$ 与 $T$，以换取更高血量和防御。第四类则是\emph{均衡型}，在总和受限时尝试避免某一维过高导致整体压缩。

值得强调的是，这种分类并不是彼此独立的自然种类，而是不同目标函数下的最优方向。例如，同样一只高 $D$ 宠，在“骑宠近战白狼”体系中可能被视作先手核心；在“骑宠远程”体系中，它的价值却可能部分让位给人物本体速度。也就是说，宠物类型并不只由宠物自身定义，而是由人物职业、武器、骑宠状态和战斗目标共同定义。

\section{人物与宠物协同}

人物与宠物协同最典型的体现，就是骑宠状态下的人宠混合公式。第 10 章已经说明：骑宠近战时，攻击和速度都更偏宠物；骑宠远程时，攻击和速度则更偏人物。因此，宠物培养的最优方向必须与人物打法协同设计。

若人物走骑宠近战，则宠物的 $S$、$T$、$D$ 都会较深地进入攻击、防御与速度路径，此时高成长宠物的收益往往能被完整兑现。若人物走骑宠远程，则宠物攻击只以较低权重并入人物攻击，而速度也更多取人物本体，因而宠物在 $D$ 维和 $S$ 维上的边际收益会相对下降。类似地，若人物更依赖职业法术控场，则宠物提供的意义可能不在纯爆发，而在于补足先手、承伤或战术容错。

从统一建模角度看，人物与宠物协同可以写成一个联合效用问题：宠物培养并不是单独最大化 $U_{pet}$，而是要最大化人宠组合的总效用
\begin{equation}
    U_{\mathrm{team}} = U_{\mathrm{char}} + U_{\mathrm{pet}} + U_{\mathrm{ride}} + U_{\mathrm{synergy}}.
\end{equation}
其中最后一项 $U_{\mathrm{synergy}}$ 正是单独看人物或单独看宠物时最容易遗漏、但在实战中最关键的部分。本章的意义，也正在于把这种协同项从“经验体感”变成可以继续分析和优化的正式对象。

